\documentclass[journal]{IEEEtran}

\usepackage{amsmath}
\usepackage{amssymb}
\usepackage{booktabs}
\usepackage{array}
\usepackage{graphicx}
\usepackage[hidelinks]{hyperref}

\newcolumntype{L}[1]{>{\raggedright\arraybackslash}p{#1}}

\begin{document}

\title{Pressure, Control, and Chaos in Synthetic Games: Construct Validity, Cross-Policy Transfer, and Falsification}

\author{Anonymous Author(s)}

\markboth{IEEE TRANSACTIONS ON GAMES}{Pressure, Control, and Chaos in Synthetic Games}

\maketitle

\begin{abstract}
We evaluate Pressure, Control, and Chaos as candidate descriptions of strategic mechanisms in four synthetic games. Historical tests involving investigator-designed policies in poker, Liar's Dice, repeated rock--paper--scissors, and spatial combat provide within-design recovery but not independent construct validation. A separate comparison uses 32 regret-minimization-derived policies and 32 logit policies without assigned construct weights. Adding three candidate scores to 95 conventional features yields no multiplicity-corrected superiority in six out-of-family predictions; for two outcomes, tests reject the hypothesis of improvement at least 0.01 training-outcome variance. Descriptive overlaps with aggression, reference regret, and entropy are substantial (absolute correlations 0.907--0.9998). An exact best-response evaluation finds 2.1533 chips-per-hand exploitability for one fixed Chaos-inspired policy snapshot, without characterizing the historical adaptive agent. These results limit claims of independent or implementation-invariant constructs. We present a falsification-oriented assessment of strategic measurements, retain unsuccessful tests, and identify requirements for stronger cross-game validation. The observed failures delimit candidate measures rather than disprove every observer-relative interpretation of strategic function.
\end{abstract}

\begin{IEEEkeywords}
Construct validity, game theory, poker, synthetic games.
\end{IEEEkeywords}

\section{Introduction}

Strategic behavior can be described through both game-theoretic quantities---best responses, regret, mixed strategies, and exploitability---and domain-specific vocabularies such as aggression, initiative, adaptation, and deception \cite{vonneumann1944,nash1950,osborne1994}. A proposed cross-game vocabulary is scientifically useful only if its definitions are sufficiently restrictive, its measurements can be evaluated independently of their labels, and its limitations remain visible when tested against established quantities.

Pressure, Control, and Chaos (PCC) originated as three candidate functional descriptions: consequential threat or commitment, context-sensitive allocation of actions, and strategically adequate unpredictability. Their motivating claim is descriptive rather than taxonomic: mechanisms might recur even where game actions and information structures differ. Yet the candidate vocabulary invites circular validation. Recovering a label that investigators used to configure a synthetic policy does not establish that an independently measurable construct exists. Similarly, a score highly correlated with aggression, regret, or entropy cannot be assumed to add explanatory or predictive content merely because it has a different name.

We therefore separate three empirical questions. \emph{Within-design recovery}: do proposed signatures respond to injected policy parameters in controlled synthetic families? \emph{Discriminant validity}: are these signatures distinct from conventional game-theoretic and behavioral variables? \emph{Out-of-family utility}: do the candidate signatures improve held-out prediction when the policies were not designed using PCC weights? These questions concern different evidence levels and cannot be answered with interchangeable statistics.

The original comparative program covered poker, Liar's Dice, repeated rock--paper--scissors (RPS), and a minimal spatial fighting game (Micro-Fighter). Its negative outcomes are retained, including inconsistent Control recovery and the absence of a tested general dominance cycle. A subsequent, separate Leduc evaluation supplies theory-agnostic CFR-derived and logit policy populations, out-of-family prediction, and exact target-conditioned best responses for one fixed Chaos snapshot. We keep those subsequent results distinct from historical policy-family claims.

This paper makes three bounded contributions. First, it states candidate PCC structures and anti-definitions while explicitly distinguishing schematic theory from estimated quantities. Second, it integrates the original four-game observations with a reproducible cross-family validation battery and equilibrium-oriented reference evaluation. Third, it documents important failures of incremental prediction, discriminant validity, and robustness, and specifies what additional evidence would be necessary before treating PCC as an independently validated cross-game construct system. We do not claim universal strategy types, human psychological validity, or a new solution concept for games.

\section{Related Work}
\label{sec:related}

The present study is positioned relative to several established lines of computational game-theory and game-AI research, rather than treating games merely as illustrative examples.

\textit{Opponent modeling and exploitation.} A substantial body of work addresses how an agent should model and respond to a specific opponent rather than compute a fixed equilibrium. Early Bayesian approaches to poker combined a prior equilibrium strategy with observations of an opponent's play \cite{southey2005}, and this was extended to large imperfect-information games through restricted Nash responses and related bounded-exploitability techniques \cite{johanson2007}, followed by related work on opponent exploitation \cite{ganzfried2011}. Safe opponent exploitation further formalizes the requirement that exploiting a fallible opponent model should not itself introduce new exploitability \cite{ganzfried2015}. This literature establishes that context-sensitive departure from equilibrium play is a well-studied, quantifiable phenomenon in poker specifically, which motivates evaluating the present Control construct through counterfactual efficiency and regret rather than through predictive accuracy alone.

\textit{Game AI and adaptive agents.} General-purpose learning systems have reached superhuman performance in Go \cite{silver2016} and, more recently, in the substantially more complex real-time domain of StarCraft II \cite{vinyals2019}. AlphaStar illustrates population-based training with specialized agents; nontransitivity and cycling are better assessed using explicit cross-play payoff matrices than inferred from strategic labels. No such complete matrix was evaluated for the proposed PCC roles here.

\textit{Behavioral and policy diversity.} Multi-agent training methods that explicitly maintain diverse populations of policies provide a related thread. Policy-space response oracles combine empirical game-theoretic analysis with deep reinforcement learning to iteratively expand a population of best-responding policies \cite{lanctot2017}, and open-ended learning methods use game-theoretic geometry to construct diverse, mutually non-dominated populations in nontransitive games such as rock--paper--scissors \cite{balduzzi2019}. These methods formalize behavioral diversity as a population-level property. The present study instead asks whether unpredictability can be a strategically productive property of a single policy---a narrower, mechanism-level question than population diversity, but one the present Chaos results suggest is compatible with it.

\textit{Procedural and synthetic environments for AI evaluation.} Synthetic and procedurally defined environments are an established methodology for evaluating general game-playing agents without overfitting evaluation to a single domain. The General Video Game AI framework, for example, provides many procedurally varied games specifically for this purpose \cite{perezliebana2019}. The present work adopts a similar comparative logic at much smaller scale: rather than many procedurally generated games, four hand-designed but structurally distinct synthetic environments are used to test whether a proposed construct generalizes, with the explicit expectation that it may not.

\textit{Bounded rationality and empirical game-theoretic analysis.} Empirical game-theoretic analysis addresses the practical difficulty of characterizing strategic interactions when the underlying game is too large or complex for direct analytical solution, instead estimating a coarser empirical game from simulated payoffs \cite{wellman2006}. Behavioral game theory similarly documents systematic ways in which real strategic behavior departs from idealized equilibrium play \cite{camerer2003}. Both literatures support treating PCC as an empirically constrained, falsifiable description of behavior rather than a claim about optimal play, and inform the present emphasis on held-out evaluation and explicit evaluation criteria without claiming a substitute for a closed-form solution concept.

\textit{Player and behavior modeling.} More broadly, player and behavior modeling is an established subfield of game AI concerned with characterizing how agents---human or synthetic---act within a game, rather than only with finding strong strategies \cite{yannakakis2018}. The present study contributes to this line of work a specific, falsifiable, cross-game comparative methodology: it does not introduce a new agent architecture or advance the state of the art in any single game, but instead asks whether three candidate descriptive mechanisms recur across games with substantially different affordances, evaluated with the same discipline---multiple policy families, retained negative controls, held-out evaluations, and separately identified fixed-opponent and best-response tests---that opponent-modeling and empirical game-theoretic research regard as relevant for credible claims.

\section{Conceptual Framework}

The candidate PCC structures are summarized as functional hypotheses about sequences, not labels attached to particular moves (Table~\ref{tab:definitions}).

\begin{table*}[t]
\caption{Candidate structural definitions and anti-definitions.}
\label{tab:definitions}
\centering
\begin{tabular}{lL{6cm}L{6cm}}
\toprule
\textbf{Construct} & \textbf{Candidate structural sequence} & \textbf{Explicit anti-definition} \\
\midrule
Pressure & Threat/commitment $\rightarrow$ response constriction $\rightarrow$ strategic consequence & Pressure $\neq$ aggression alone \\
Control & Information $\rightarrow$ context $\rightarrow$ allocation among strategically viable responses & Control $\neq$ prediction or retreat alone \\
Chaos & Unpredictability + strategic adequacy + resistance to exploitation & Chaos $\neq$ randomness \\
\bottomrule
\end{tabular}
\end{table*}

\subsection{Pressure}

Pressure concerns the alteration of an opponent's strategic environment through credible threat or commitment. A candidate Pressure process contains three stages: exposure to threat or commitment ($T$), constriction or redistribution of viable responses ($R$), and strategic consequence ($V$). The requirement for consequence is essential. Frequent attacks, large bets, or forward movement do not automatically constitute Pressure if they do not change the opponent's strategically viable choices or the costs associated with those choices.

Response constriction need not mean a literal reduction in legal actions. In poker, the legal menu may remain fold, call, or raise while commitment changes their relative costs and feasibility. In Liar's Dice, escalation changes the attractiveness of continuing versus challenging. In spatial combat, threat can compress movement and defensive options. Thus the structural hypothesis is game-independent only at a high level; the observable mechanism remains game-specific:
\begin{equation}
P: T \rightarrow R \rightarrow V \quad \text{(schematic only)}
\end{equation}

\subsection{Control}

Control concerns the context-sensitive use of strategically relevant information. The candidate sequence is information uptake ($I$), contextual alignment ($K$), and allocation among strategically viable responses ($A$). The final stage is intentionally phrased as allocation rather than intervention. Early experiments emphasized prediction and active counteraction, but later poker and Micro-Fighter results showed that a contextually correct Control response may be passive: checking, calling, delaying commitment, or preserving initiative can be more valuable than acting aggressively.

This formulation distinguishes Control from prediction. An agent can forecast an opponent accurately yet exhibit little Control if those forecasts do not materially alter action allocation. It also distinguishes Control from a fixed defensive heuristic. Retreating under pressure, for example, can reduce value if the retreat is non-persistent or forfeits initiative. Control therefore requires both context sensitivity and strategic appropriateness:
\begin{equation}
C: I \rightarrow K \rightarrow A_{\mathrm{alloc}} \quad \text{(proposed causal sequence).}
\end{equation}

\subsection{Chaos}

Chaos requires the strongest anti-definition: Chaos is not randomness. Shannon entropy provides a natural measure of uncertainty in a distribution \cite{shannon1948}, but entropy alone does not establish strategic function. A uniformly random policy may be maximally unpredictable while being systematically weak or easily exploitable in an environment where many random actions are poor.

The emerging Chaos construct therefore contains three requirements: behavioral unpredictability ($U$), strategic adequacy ($S$), and resistance to exploitation ($E$). These dimensions can be written schematically as
\begin{equation}
Ch = f(U, S, E),
\end{equation}
where $f$ is an unestimated candidate aggregation function, not a fitted model. A multiplicative form is a hypothesis, not an established result; additive and multiplicative specifications require prospective comparison on held-out data. This conception is related to, but distinct from, equilibrium randomization. Mixed strategies are foundational in game theory \cite{vonneumann1944,nash1950}, yet an equilibrium mixture need not reveal a latent Chaos mechanism, and a strategically effective Chaos policy need not maximize raw entropy.

\subsection{Non-universality and evidence levels}

PCC does not require all three constructs to exist in every environment. Repeated RPS illustrates why this matters: under the current formulation there is no meaningful analogue of escalating commitment or persistent threat sufficient to instantiate Pressure, so Pressure is treated as absent by design rather than manufactured through an arbitrary observable.

The program also distinguishes mechanism evidence from construct recovery. A causal intervention can demonstrate that a candidate mechanism behaves as predicted without establishing that the construct can be inferred robustly from observational behavior. Stronger evidence requires recovery across independently structured policy families; stronger still is cross-game portability. Human validity requires separate external evidence. This hierarchy prevents successful synthetic demonstrations from being silently promoted into claims they do not support.

\section{Comparative Methods}
\label{sec:methods}

The comparative PCC program was designed around the risk of post-hoc construct fitting. Candidate mechanisms were implemented in synthetic policy families, measured through public or policy-independent observables where practical, evaluated using held-out seeds, and evaluated using historically recorded acceptance criteria whose dated provenance is incomplete. Failed, partial, unresolved, and absent-by-design outcomes were retained rather than repaired automatically through threshold adjustment. When a result failed, subsequent work could diagnose the failure, but follow-up diagnostics and interventions were recorded separately; their original confirmatory status cannot be verified from the available timestamp evidence.

The four environments were selected to span distinct strategic affordances (Table~\ref{tab:environments}). Poker combines imperfect information, stochastic state transitions, sequential betting, commitment, opponent adaptation, and mixed strategy. Liar's Dice preserves bluffing, hidden information, escalation, and challenge timing while removing poker-specific streets, ranges, and pot construction. Repeated RPS provides a deliberately minimal repeated game and a natural test of mixed behavior and learning; experimental work on matching-pennies-type games shows that observed behavior can depend materially on information about past opponent actions \cite{mookherjee1994}. Micro-Fighter adds a qualitatively different substrate---space, range, cooldown, health, initiative, defense, and counterattack---without labeling any move as intrinsically Pressure, Control, or Chaos.

\begin{table*}[t]
\caption{Experimental environments and their roles.}
\label{tab:environments}
\centering
\begin{tabular}{lL{4.2cm}L{8.5cm}}
\toprule
\textbf{Environment} & \textbf{Primary methodological role} & \textbf{Key affordances} \\
\midrule
Poker & Deep construct-recovery laboratory & Imperfect information, stochasticity, sequential commitment, opponent adaptation \\
Liar's Dice & Near-domain transfer & Hidden dice, escalating public claims, bluffing, challenge timing \\
Repeated RPS & Negative control / identifiability & Minimal repeated action space, history dependence, high-entropy baseline \\
Micro-Fighter & Cross-domain spatial probe & Distance, threat range, cooldown, initiative, defense, counterplay \\
\bottomrule
\end{tabular}
\end{table*}

\subsection{Poker}

Poker served as the deepest laboratory. The synthetic program used multiple independently structured policy families and supported mixture recovery, temporal-history manipulations, context-yoked models, matched counterfactual interventions, fixed-opponent exploitation tests, and value-sensitive analyses. Poker is a canonical imperfect-information domain in computational game theory: counterfactual regret minimization was developed and demonstrated at large scale in poker \cite{zinkevich2007}, heads-up limit hold'em was later essentially solved \cite{bowling2015}, and superhuman systems were subsequently developed for heads-up no-limit and multiplayer no-limit poker \cite{brown2018,brown2019}. The present PCC work does not attempt to solve poker; it uses poker as a structured environment in which latent synthetic mechanisms can be manipulated and evaluated.

\subsection{Liar's Dice}

Liar's Dice was used as a near-domain transfer environment. Players possess hidden dice, make escalating public claims, and choose whether to continue bidding or challenge. Two independently structured synthetic families were evaluated for competitiveness, construct recovery, Control--Chaos decomposition, and Pressure-related response constriction. The purpose was not to reproduce poker mechanics but to test whether structurally related mechanisms survive when commitment is expressed through claims rather than monetary betting.

\subsection{Repeated RPS}

Repeated RPS was designed primarily as a falsification and identifiability environment. The action space is minimal, there is no persistent resource state, and Pressure was declared absent by design. RPS was especially useful for testing Chaos because uniform random play provides a natural high-entropy baseline. The environment asks whether a proposed Chaos measure can distinguish strategically functional unpredictability from ordinary randomization or equilibrium-like mixing.

\subsection{Micro-Fighter}

Micro-Fighter was developed as a headless, deterministic-tick, minimal spatial combat environment with identical agents and a shared action set (advance, retreat, attack, guard, evade). State variables include distance, health, cooldowns, and public action history. It enables tests of spatial constriction, threat generation, defensive counter-value, initiative preservation, and adaptive exploitation. This environment was deliberately chosen to test whether PCC-like mechanisms survive outside card and dice settings.

\subsection{Historical evaluation workflow}

Major experiments followed a design $\rightarrow$ freeze criteria $\rightarrow$ evaluate $\rightarrow$ retain outcome workflow. Historical protocols describe acceptance criteria and seed separation for selected synthetic runs; complete dated freeze artifacts and their provenance are not established in this manuscript, so their confirmatory status cannot currently be independently audited. Reported failures were retained in the archival comparative record. Later decompositions investigated possible mechanisms; the archived descriptions report unchanged evaluation gates, but the original dates and thresholds are not fully independently verifiable. This procedure was especially important in the Poker Control and Micro-Fighter Control sequences.

\subsection{Latent mechanisms and observables}

Synthetic policies expose their latent mechanism weights for validation, but these assignments were not used directly as measurements. Candidate observables were derived from public actions, public state, model predictions, payoffs, or matched counterfactuals. Latent weights were used only after aggregation to assess construct validity. This separation prevents direct label leakage into the observable, but does not establish independent construct validity: correlations with researcher-assigned latent weights are within-design parameter recovery, not analyst-blind recovery of a naturally occurring construct.

\subsection{Held-out evaluation and exploiters}

Calibration and evaluation seeds were separated whenever an experiment involved model fitting, candidate selection, or adaptive exploitation. Exploiters were calibrated only against designated predictable baselines and frozen before evaluation against Chaos-oriented policies. This is a fixed-opponent stress test, not a best-response exploitability estimate. Since this exploiter was not trained against each target policy, comparisons cannot establish worst-case robustness.

Historical reports distinguish within-design support, partial support, failure, and untested conditions. Retrospective threshold provenance is incomplete, and these labels do not substitute for inferential uncertainty.

\subsection{Interpretive safeguards and outstanding validation}
The original correlation coefficients compare observed signatures with latent parameters assigned by the investigators. They are therefore \emph{self-recovery} coefficients; neither analyst blindness nor ecological construct validity follows from such comparisons. Their inferential interpretation is limited further by the absence of documented independent replication units, seed counts, confidence intervals and familywise multiplicity adjustments in the present archival summary. Historical ``pass'' labels are reported as original workflow outcomes, not statistically established acceptances. Similarly, claims of preregistration cannot be audited until immutable, dated threshold files and exact evaluation configurations are supplied. The coefficient 3.35 appears in a follow-up intervention after failures had been observed; no derivation or timestamp proving advance selection is available in the archived record. We classify this intervention and the subsequent revisions as exploratory.

The fixed-opponent payoff tests do not measure worst-case exploitability. Target-specific best responses, approximate-equilibrium opponents and matched-core perturbation controls are outstanding. Finally, the manuscript reports no measured time series of $[P(t),C(t),Ch(t)]$ and no dominance-cycle meta-game matrix. Neither dynamic tripartite emergence nor a universal dominance cycle is established by the current evidence.

\subsection{Separate equilibrium-reference infrastructure}
\label{sec:newcfr}
A newer poker research snapshot, not used to generate the original cross-game result tables, contains a two-player limit Leduc counterfactual-regret-minimization (CFR) benchmark. The six-card deck contains two cards of each of three ranks; the implementation uses two betting rounds with bet/raise increments of 2 and 4 chips, respectively, and permits at most two raises after an initial bet within a round (three bets/raises total, one more than the common two-bet variant). Its saved convergence artifact describes 90 chance deals and 528 unique information sets. The internally computed average-policy exploitability (defined in the artifact as NashConv/2, in net chips per hand) decreases from 3.1306 at one iteration to 0.2161 at 100 and 0.06977 at 500 iterations. These are \emph{solver profile} measurements, not target-conditioned best-response exploitabilities of the historical PCC Chaos policies. The saved report itself states that external solver cross-validation has not been performed. Thus the newer artifact addresses the absence of a concrete equilibrium-oriented reference implementation only partially; it cannot retroactively validate the original frozen-exploiter comparisons or other game environments. The artifact is listed with its SHA-256 checksum in the accompanying evidence manifest.

\subsection{Prospective cross-family prediction and fixed-policy best response}

A separate experiment recorded and froze its measurement code, seeds, ridge penalty, endpoints, and a 12-test comparison family before generating 32 chance-sampled CFR policies and 32 independently specified equity/pot-odds logit policies. Neither family had assigned PCC weights. A sealed identity map was read only after measurements and predictions were hashed. This was computational label withholding, not independent investigator blindness: the implementing analyst knew both generators, and cohort identity could be inferred.

The analysis trained on one entire family and tested on the other, in both directions. A public-history split assigned 300 information sets to candidate measurements and 228 to a held-out matched-card response endpoint. The conventional baseline used 95 features covering aggression, entropy, action frequencies, history strata, reference regret, and exact exploitability. Three new reference-based candidate PCC scores measured wager commitment/fold opportunity, uniform-reference decision efficiency, and adequacy-weighted surprisal. They operationalize restricted candidates, not the full theoretical constructs. Outcomes were exact seat-balanced value against separately fixed tight and loose logit opponents, and held-out matched-card total variation. Ridge fitting and standardization used only the training family. Whole-policy bootstrap resampling independently resampled training and testing populations and refit the paired models; hands and information sets were not independent replicates.

Define improvement as baseline minus augmented test squared error divided by training outcome variance. Positive values favor adding the candidate scores. Table~\ref{tab:prospective-transfer} reports all six estimates and marginal percentile bootstrap intervals from 3,999 resamples. No superiority test rejected after Holm correction over six superiority and six minimum-benefit tests (all adjusted superiority $p=1.000$). The one-sided minimum-benefit null was improvement at least 0.01, one percent of training outcome variance; it was rejected only for the two logit-to-CFR value outcomes (adjusted $p=0.006$ and $0.0165$). The 95\% intervals are uncentered percentile quantiles of 3,999 paired two-population bootstrap refits (2.5th and 97.5th percentiles). One-sided $p$-values instead use the centered bootstrap distribution $b^*-\hat\Delta$: for the minimum-benefit null $\Delta\geq0.01$, they count centered draws not exceeding $\hat\Delta-0.01$, with a plus-one correction, then apply Holm across twelve tests. These are different inferential constructions; neither a marginal two-sided interval nor its upper endpoint is an inversion of the Holm-adjusted one-sided test. Thus these results do not establish incremental construct validity. Failure to reject superiority does not prove exactly zero benefit, and transfer in the reverse direction remains imprecisely estimated.

\begin{table}[t]
\caption{New cross-family incremental prediction. Each direction has 32 training and 32 testing policies. Intervals are marginal 95\% paired population-bootstrap intervals; multiplicity-adjusted tests are described in the text.}
\label{tab:prospective-transfer}
\centering\scriptsize
\setlength{\tabcolsep}{3pt}
\begin{tabular}{llrr}
\toprule
Train $\to$ test & Outcome & Improvement & 95\% interval \\
\midrule
CFR $\to$ logit & Tight value & $-0.3867$ & $[-6.2615,0.4625]$ \\
CFR $\to$ logit & Loose value & $0.1032$ & $[-2.8740,1.5138]$ \\
CFR $\to$ logit & Card response & $-2.4863$ & $[-10.4523,2.2479]$ \\
Logit $\to$ CFR & Tight value & $-0.0235$ & $[-0.0341,0.0189]$ \\
Logit $\to$ CFR & Loose value & $-0.0311$ & $[-0.0530,0.0236]$ \\
Logit $\to$ CFR & Card response & $-0.0087$ & $[-0.0232,0.0120]$ \\
\bottomrule
\end{tabular}
\end{table}

Separately, a fixed behavioral snapshot of the source Score Chaos policy, with mixture $(0.1,0.1,0.8)$, temperature 0.35, and frozen empty opponent and action histories, was evaluated directly in the solver's Leduc game. This snapshot is not the history-updating policy from the historical tournaments. Its profile value for player 0 was $-0.05255$ net chips per hand. Exact target-conditioned best-response values for players 0 and 1 were 1.96121 and 2.34530; their sum gives NashConv 4.30651 and exploitability 2.15326. These are deterministic full-enumeration quantities conditional on the specified policy and rules, not Monte Carlo estimates requiring hand-level confidence intervals.

A separate recursive response algorithm reproduced both values. An independent dictionary-state rule evaluator checked all 14,490 solver nodes, information keys, legal actions, terminal payoffs, and chance probabilities, including agreement between 120 physical-card chance triples and 90 rank-deck arrangements. Executing both response tables in that evaluator agreed within $10^{-9}$ chips. Eleven focused tests passed for label withholding, feature isolation, dependence, measurement invariance, and rule/response consistency. This is internal cross-implementation checking, not external OpenSpiel replication. Sixty-five pre-existing validation files retained their hashes. The new experiment neither changes historical results nor establishes analyst-blind validation, full feedback Control, or cross-game generalization.

\subsection{Post-reveal measurement overlap and incremental utility}
\label{sec:overlap}
The 64-policy prospective study also produced descriptive, pooled post-reveal correlations between the three candidate PCC scores and conventional measurements. Candidate Pressure was strongly correlated with aggression ($r=0.99983$), candidate Chaos with entropy ($r=0.99090$), and candidate Control was inversely correlated with reference regret ($r=-0.90684$). These are descriptive pooled correlations, not multiplicity-adjusted confirmatory tests; their direction and magnitude depend on the evaluated policy population and candidate score definitions. They nevertheless undermine any assertion that these particular scores are independently established constructs. They also explain why favorable within-family self-recovery would be insufficient to establish distinctiveness.

\begin{table}[t]
\caption{Descriptive candidate-to-baseline measurement overlap in the post-reveal 64-policy pool. These values were not prospectively tested for statistical significance.}
\label{tab:overlap}
\centering\small
\begin{tabular}{llr}
\toprule
PCC candidate & Conventional feature & Correlation \\
\midrule
Pressure & Aggression & $0.99983$ \\
Control & Reference regret & $-0.90684$ \\
Chaos & Entropy & $0.99090$ \\
\bottomrule
\end{tabular}
\end{table}

\section{Results}

The four environments produced markedly different recovery profiles, as summarized in Table~\ref{tab:crossgame}. No environment supplied implementation-invariant validation of all three constructs. A Pressure--Control--Chaos dominance cycle was not directly tested because no complete head-to-head cross-play matrix was constructed. This non-uniformity is not treated as a defect to be removed; it is a primary result. The strongest commonalities occur at the structural level rather than at the level of one shared scalar observable.

\begin{table*}[t]
\caption{Conservative cross-game PCC status.}
\label{tab:crossgame}
\centering
\begin{tabular}{lL{3cm}L{5cm}L{5cm}}
\toprule
\textbf{Environment} & \textbf{Pressure} & \textbf{Control} & \textbf{Chaos} \\
\midrule
Poker & Reported within-design support & Partial / implementation-dependent & Exploratory fixed-opponent evidence \\
Liar's Dice & Partial & Failed & Reported within-design support \\
Repeated RPS & Absent-by-design & Failed & Not separately identifiable from equilibrium play \\
Micro-Fighter & Exploratory mechanism evidence & Exploratory mechanism evidence & Exploratory mechanism evidence \\
\bottomrule
\end{tabular}
\end{table*}

\subsection{Poker}

\subsubsection{Pressure}

Pressure was the strongest and most stable poker construct. Historically reported synthetic validation recovered Pressure-related observables across independent implementations, and mechanistic experiments supported the relationship among commitment, opponent response, and strategic consequence. Later investigations did not require revision of the Pressure definition, so Pressure has comparatively strong within-design mechanism evidence. However, the newer candidate score is almost perfectly correlated with aggression in the independent-policy pool; distinct construct validity has not been established.

\subsubsection{Control}\label{sec:poker-control}

Control produced the most complicated sequence. A structural-recovery experiment decomposed Control into information uptake, contextual alignment, and value-sensitive allocation. The Adaptive family passed all three stages, with approximate correlations to latent Control weight of 0.654, 0.543, and 0.381. The independently structured Score family initially failed all three, with corresponding correlations of approximately 0.003, $-0.023$, and $-0.087$. This family split suggested that the failure was not simply a weak universal measurement.

A direct architectural decomposition held poker state fixed while varying learned opponent fold probability from 0.10 to 0.90. Mean total-variation shift in Score-Control's action distribution was approximately 0.164, compared with approximately 0.998 for Adaptive-Control. The observed shift ratio was approximately 6.1; the near-ceiling Adaptive-Control shift could also signal brittleness, which requires perturbation testing rather than an automatic favorable interpretation. This localized the problem to policy architecture: Score-Control remained dominated by card/value terms, whereas Adaptive-Control allowed learned response information to reshape action selection.

An exploratory contextual intervention added the zero-centered term $\Delta = 3.35(p_{\text{fold}} - 1/3)$ to Score-Control aggression while leaving the remaining value/risk terms intact. After the change, Score-Control was reported to pass information uptake ($r \approx 0.358$) and contextual alignment ($r \approx 0.315$) but still failed value-sensitive allocation ($r \approx 0.132$) at this first checkpoint; Table~\ref{tab:pokercontrol} reports the lower final values retained after the two further exploratory tests described below. A value decomposition then showed that Control weight was negatively associated with counterfactual efficiency ($r \approx -0.754$) and positively associated with regret ($r \approx +0.533$), even while positive contextual gain increased with Control ($r \approx +0.220$). These negative value associations raise the alternative possibility that the parameter designated Control increased poor decisions or that the operational label was misassigned; they do not independently establish successful Control.

Two subsequent exploratory tests used the original recovery gate. Hard-gating contextual aggression by a frozen efficiency threshold of 0.80 did not resolve Control; information uptake was approximately 0.196, context alignment 0.256, and value-sensitive allocation 0.101. A matched-state comparison then showed that Adaptive-Control achieved approximately $+0.460$ greater one-step value on the same states and preserved roughly $+0.170$ additional check/call probability against fold-resistant opponents. This suggested that Adaptive-Control used context on both sides of the decision---timed aggression against vulnerable opponents and preserved optionality against resistant ones. A final two-sided Score intervention retained the aggression term and added passive optionality, but the final value-sensitive stage remained below the frozen criterion ($r \approx 0.141$). No further coefficient tuning was performed.

The conservative conclusion is architecture-dependent within-design recovery, with an important unresolved contradiction between Control weight and value. The present experiments do not establish implementation-invariant Control or distinguish a construct-measurement failure from a defective policy parameterization.

\subsubsection{Chaos}

Poker Chaos supplied a separate fixed-opponent comparison. A falsification test compared predictable value-oriented play, uniform random play, Score-family Chaos, and an independently structured Chaos policy. Uniform random play had normalized entropy 1.000 but mean payoff $-1.447$ against a neutral baseline and $-1.794$ against a frozen exploiter. Score Chaos had lower entropy (0.908) but substantially better payoffs ($-0.300$ and $-0.439$). Independent Chaos had still lower entropy (0.539) yet positive payoffs ($+0.432$ and $+0.737$). The exploiter was calibrated only against predictable play and frozen before Chaos evaluation.

Both Chaos families satisfied the originally reported fixed-opponent screening conditions, interpreted here as exploratory descriptive gates rather than proof of exploitability resistance. These comparisons demonstrate that entropy alone does not determine payoff in these particular matchups. They cannot establish low exploitability: the designated exploiter was frozen after calibration on other policies and is not a target-specific best response. In particular, the positive payoff advantage of Independent Chaos against this exploiter could reflect opponent mismatch or overall policy strength. Matched-value-core experiments and best-response analysis are needed.

\begin{table*}[t]
\caption{Poker Control structural-recovery summary. Score-family values are the final retained figures after all three exploratory intervention stages, not the first checkpoint described in Section~\ref{sec:poker-control}.}
\label{tab:pokercontrol}
\centering
\begin{tabular}{lcc}
\toprule
\textbf{Poker Control stage} & \textbf{Adaptive family} & \textbf{Score family (final retained)} \\
\midrule
Information uptake & Passed ($r \approx 0.654$) & Passed ($r \approx 0.307$) \\
Context alignment & Passed ($r \approx 0.543$) & Passed ($r \approx 0.220$) \\
Value-sensitive allocation & Passed ($r \approx 0.381$) & Failed ($r \approx 0.141$) \\
\bottomrule
\end{tabular}
\end{table*}

\begin{table*}[t]
\caption{Poker Chaos falsification: randomness versus effective unpredictability.}
\label{tab:pokerchaos}
\centering
\begin{tabular}{lccc}
\toprule
\textbf{Policy} & \textbf{Normalized entropy} & \textbf{Payoff vs.\ neutral} & \textbf{Payoff vs.\ frozen exploiter} \\
\midrule
Predictable value & 0.000 & $-0.659$ & $-0.716$ \\
Uniform random & 1.000 & $-1.447$ & $-1.794$ \\
Score Chaos & 0.908 & $-0.300$ & $-0.439$ \\
Independent Chaos & 0.539 & $+0.432$ & $+0.737$ \\
\bottomrule
\end{tabular}
\end{table*}

\subsection{Liar's Dice}

Liar's Dice produced a mixed profile: Pressure partial, Control failed, and reported within-design support for Chaos. Pressure-like mechanisms appeared through escalating claims and response effects, but response constriction was implementation-sensitive rather than invariant across both policy families. The tested Control operationalization did not satisfy full recovery requirements and was retained as a failure rather than redefined after the fact. For the tested agents, bounded unpredictability produced more favorable payoffs against a designated fixed opponent than the reported naive stochastic comparator. This is an additional within-design imperfect-information example, not independent confirmation of a general Chaos construct or low worst-case exploitability.

\subsection{Repeated Rock--Paper--Scissors}

Repeated RPS yielded the most negative PCC profile: Pressure absent-by-design, Control failed, and Chaos failed as a fully identifiable latent construct. This negative profile is methodologically valuable. Pressure was not manufactured from arbitrary action statistics in an environment lacking escalating commitment or persistent threat. History-sensitive behavior could exploit predictable opponents, but adaptive prediction alone did not satisfy the Control construct. For Chaos, uniform random play supplied a key counterexample: maximally unpredictable behavior can be strategically adequate in the game-theoretic sense while still revealing little about a latent Chaos mechanism. In the one-shot zero-sum RPS stage game, uniform mixing is the unique Nash equilibrium. Its strategic adequacy is therefore expected, and the inability to infer a distinct latent Chaos mechanism from uniform play is an observational-identifiability issue rather than a failure of equilibrium play.

\subsection{Micro-Fighter}

Micro-Fighter tested PCC in a spatial domain. Pressure-oriented behavior generated spatial compression, in-range attack opportunities, and forced defensive responses, supporting the early stages of the Pressure structure. However, threat volume did not universally convert into damage or victory: one Control implementation could absorb spatial pressure and counter effectively. Thus spatial aggression is insufficient without strategic consequence.

Control generated an especially instructive falsification. An exploratory retreat-under-sustained-pressure intervention worsened performance. Approximately 80\% of retreat triggers failed to increase distance; after successful distance gain, Pressure re-entered immediately in roughly 88\% of cases; only about 1\% of successful retreats preserved separation for two ticks; and roughly 88\% of retreat triggers plausibly forfeited an attack opportunity. These outcomes are compatible with a kinematic or hand-coded-policy artifact. Without movement-speed, arena-geometry and opponent-sensitivity ablations, they establish only a limitation of this implementation, not of spatial Control generally.

Chaos produced exploratory mechanism-level evidence. A state-aware random fighter had higher normalized conditional entropy (about 0.933) than an effective-Chaos candidate (about 0.878), yet decisively underperformed against a neutral opponent (about 0.031 vs.\ 0.418 decisive win rate). A stronger adaptive exploiter, calibrated only against predictable play, yielded held-out decisive win rates of 0.000 for predictable play, 0.381 for random play, and 0.720 for effective Chaos; effective Chaos also maintained a positive health margin of approximately $+1.015$. Full cross-family recovery and target-specific exploitability remain untested. These fixed-opponent results are compatible with the effective-Chaos policy simply having a stronger underlying value policy.

\section{Cross-Game Synthesis}
\label{sec:synthesis}

Taken together, the four environments do not support a universal dominance relation. The strongest cross-game regularities instead appear as structural motifs. Pressure is best represented as threat or commitment followed by response constriction and strategic consequence. Control is best represented as information uptake followed by contextual interpretation and allocation among viable responses. Chaos is best represented as unpredictability coupled to value preservation and resistance to exploitation.

The absence of one shared scalar is not itself evidence against portability. The same strategic structure can be realized through different observables when the game changes. What matters is whether the causal or functional sequence survives. At the same time, this flexibility creates the risk of post-hoc equivalence claims. The framework therefore treats cross-game similarity as provisional unless the mapping is prespecified and independently tested.

In poker and Micro-Fighter, greater policy entropy does not guarantee greater performance against the particular evaluated opponents. The poker uniform policy has entropy 1.000 but loses more to the tested opponents than an independently designed policy of entropy 0.539 (Table~\ref{tab:pokerchaos}); the fixed opponent is not a best-response oracle. In Micro-Fighter, the random policy has conditional entropy 0.933 and a decisive win rate of 0.031 against neutral, versus entropy 0.878 and a win rate of 0.418 for the effective-Chaos candidate. Against the calibrated exploiter the respective win rates are 0.381 and 0.720: the gap \emph{narrows} from 0.387 to 0.339, and both policies improve against this opponent. These comparisons confound stochasticity with policy quality and do not establish target-independent exploitability resistance.

Uniform mixing in standard rock--paper--scissors is a Nash equilibrium, not a strategic failure. Results from this game therefore cannot support the claim that higher entropy universally reduces value. Across the four settings, the justified conclusion is narrower: entropy alone does not operationalize a distinct Chaos mechanism. Likewise, without a full pairwise PCC-policy payoff matrix, a universal dominance cycle is not tested and cannot be said to fail cross-game transfer.

Table~\ref{tab:falsifications} summarizes the tested counterexamples and untested boundary claims. It is a map of evidentiary limitations rather than proof of a new strategic taxonomy.

\begin{table*}[t]
\caption{Principal falsifications and boundary conditions.}
\label{tab:falsifications}
\centering
\begin{tabular}{L{5.5cm}L{9.5cm}}
\toprule
\textbf{Falsified or rejected simplification} & \textbf{Comparative finding} \\
\midrule
Chaos = randomness & High entropy can coexist with poor payoff against fixed opponents; high exploitability was measured only for a separate frozen Leduc Chaos snapshot. \\
Control = prediction & Prediction without value-sensitive allocation does not recover Control. \\
Control = retreat & Spatial retreat can forfeit initiative and worsen outcomes. \\
Pressure = aggression & Threat generation without response constriction or consequence is insufficient; the tested candidate Pressure score correlates with aggression at $r=0.99983$. \\
All games contain all three constructs & Repeated RPS provides an absent-by-design Pressure case. \\
PCC is a universal dominance cycle & No complete cross-play matrix was evaluated; the claim remains untested. \\
\bottomrule
\end{tabular}
\end{table*}

\section{Discussion}

\subsection{PCC as mechanisms rather than strategy types}

The historical findings are compatible with viewing PCC as candidate functional descriptions rather than fixed player types; they do not validate an independent three-dimensional decomposition. A single policy may exhibit commitment, context-sensitive allocation, and randomized action selection, but this observation alone does not establish three distinct latent constructs. A time-varying PCC vector remains a theoretical proposal: this manuscript estimates no such trajectory.

A fixed cyclic dominance relation cannot be evaluated without a complete cross-play payoff matrix. The relative value of constraint, contextual adaptation, and unpredictability depends on the affordances of the game. A construct can be present without enjoying a universal advantage over another construct.

\subsection{Independent construct validity and conventional explanations}
\label{sec:incremental}
The central new empirical outcome is not the recovery of PCC labels but the failure of the restricted candidate score bundle to demonstrate incremental predictive superiority over a substantial conventional feature set in the two-direction cross-family experiment (Table~\ref{tab:prospective-transfer}). Absence of corrected significance is not evidence for exact equivalence. Yet the post-reveal score overlap (Table~\ref{tab:overlap}) provides an additional, conceptually independent warning: these implementations closely track standard aggression, reference regret, and entropy features. Because the baseline already contains game-theoretic measurements with direct operational meanings, a convincing claim of independent PCC utility would require prespecified out-of-family predictive gains, distinct intervention predictions, or improved scientific explanation not furnished by those measurements. None is established here.

The frozen Score Chaos snapshot likewise illustrates why action diversity and resistance to exploitation must not be conflated. Its exact exploitability is a property of a particular fixed snapshot under specified Leduc rules, not a measurement of all Chaos-inspired policies or of their adaptive versions. The historically reported frozen-opponent payoff contrasts remain opponent-specific rather than general exploitability estimates. Taken together these findings make conditional and negative evidence central to the contribution, rather than exceptions to an assumed successful taxonomy.

Failures of Control recovery, naive randomness, and strategic retreat narrow the supported claims, but do not eliminate researcher degrees of freedom. Historical thresholds without independently available freeze records are not described as prospectively validated.

\subsection{Control as the hardest construct}

Control appears to be the most implementation-dependent PCC component. Its difficulty may be intrinsic to observation. Threat can often be measured through commitment or spatial exposure, and unpredictability can be quantified through behavioral distributions and exploitation. Control, by contrast, often requires the counterfactual question: given this information and this context, did the agent allocate behavior differently in a way that was appropriate relative to other viable actions? That is a harder inference problem, and the poker sequence in Section~\ref{sec:poker-control} shows why it cannot be dismissed as a simple failure to encode opponent information: contextual sensitivity was successfully induced without resolving value-sensitive allocation, pointing toward a richer concept of strategic allocation that may require modeling the alternative actions a policy declines to take.

The availability of a proposed mechanism depends on game affordances, and labels cannot be assigned simply to complete a three-part taxonomy.

\subsection{Relation to established game-theoretic ideas}

PCC should not be read as a replacement for equilibrium analysis. Classical game theory already provides rigorous concepts of mixed strategy, best response, equilibrium, regret, and exploitation \cite{vonneumann1944,nash1950,osborne1994}. Computational poker research further demonstrates how counterfactual regret and exploitability can be formalized at enormous scale \cite{zinkevich2007,bowling2015,brown2018,brown2019}. PCC operates at a different descriptive level: it asks whether recurring functional mechanisms can organize behavior across games without assuming that those mechanisms correspond one-to-one with equilibrium concepts.

The Chaos construct is the closest to mixed-strategy reasoning, but the difference is important. Randomization can be equilibrium-relevant without implying a recoverable Chaos mechanism. Conversely, structured nonuniform variation can be strategically hard to exploit without maximizing entropy. The Control construct overlaps with learning and best-response ideas, but the present results emphasize context-sensitive allocation rather than prediction alone. Pressure overlaps with commitment and threat, yet the present comparative criterion additionally requires evidence of response constriction and consequence.

All agents evaluated here are synthetic. The results neither validate human psychological constructs nor generalize to human gameplay.

\section{Limitations}

Several limitations constrain interpretation. First, the synthetic policy families were engineered within a PCC research program, so even independently structured implementations are not theory-neutral. Second, the number of environments is small. Third, cross-game structural mappings such as ``response constriction'' are conceptually aligned but not numerically identical. Fourth, Control remains incompletely implementation-invariant. Fifth, Micro-Fighter is a deliberately minimal synthetic environment rather than an established competitive game. Sixth, adaptive exploiters depend on their model class and calibration design. Seventh, latent strategic intent may be fundamentally non-identifiable when different mechanisms induce equivalent behavior. Finally, the present paper contains no human behavioral evidence supporting the cross-game claims. 

The paper therefore does not establish PCC as a universal law of competition or as three empirically separable latent dimensions. The defensible result is narrower: PCC-inspired hypotheses can be studied in several synthetic settings, but candidate measurements exhibit architecture dependence, substantial conventional-feature overlap, and unsuccessful out-of-family incremental prediction in the evaluated populations.

\section{Falsification Criteria and Future Tests}

Confidence in the framework should decrease if Pressure-like behavior repeatedly fails to constrain opponents or produce strategic consequences; if Control-like behavior is recoverable equally well without relevant information or context; if strategically useless random behavior consistently satisfies Chaos criteria; if the structural definitions require incompatible meanings in each new domain; if independent implementations systematically fail construct recovery; or if prospectively specified predictions perform no better than arbitrary post-hoc descriptions.

Further tests should use prospectively frozen designs. A new competitive environment should be selected in advance, PCC predictions should be specified before construct-specific policies or observables are tuned, and the predictions should be evaluated under independent replication. PCC becomes a more credible theory only to the extent that it predicts strategic structure before that structure is observed.

\section{Conclusion}

This comparative case study investigated whether three proposed descriptions of strategic mechanism---Pressure, Control, and Chaos---could survive increasingly independent validation. The historical synthetic experiments yielded conditional self-recovery and several retained failures, but did not provide implementation-invariant construct validation. A separate study evaluated 64 policies not assigned PCC weights: six cross-family comparisons showed no corrected incremental superiority over 95 conventional predictors. Two tests rejected the hypothesis that improvement was at least the prespecified practical-benefit threshold, while post-reveal score overlap indicated near-redundancy for some operationalizations. Exact best responses also found substantial exploitability for a separately fixed Chaos policy snapshot, without warranting conclusions about the historical adaptive policy.

The result is neither proof that all observer-relative interpretations of PCC are false nor evidence that they are independently validated. Rather, it establishes concrete limits of the current game-level measurements and demonstrates a workflow in which negative results constrain a proposed strategic vocabulary. Future claims of distinctiveness require independently specified construct estimators, provenance-complete historical protocols, genuinely external benchmarks, and tests of prospectively defined incremental predictions beyond established game-theoretic variables. Until then, PCC in these synthetic games is best treated as a candidate comparative vocabulary under test, not a validated alternative to equilibrium or opponent-modeling theory.

\section*{Data and Code Availability}
The accompanying supplementary package includes the available synthetic-policy validation artifacts, protocol and source manifests, and analysis outputs for the newer policy-transfer study. The original four-game historical study is not yet fully reproducible from a single version-pinned archive: some dated acceptance thresholds, seeds, and policy configurations could not be independently reconstructed. No human gameplay data are analyzed. The supplementary package documents these limitations rather than claiming complete reproduction of the historical experiments.

\begin{thebibliography}{23}

\bibitem{vonneumann1944}
J. von Neumann and O. Morgenstern, \emph{Theory of Games and Economic Behavior}. Princeton, NJ, USA: Princeton Univ. Press, 1944.

\bibitem{nash1950}
J. F. Nash, ``Equilibrium points in $n$-person games,'' \emph{Proc. Nat. Acad. Sci.}, vol. 36, no. 1, pp. 48--49, 1950.

\bibitem{osborne1994}
M. J. Osborne and A. Rubinstein, \emph{A Course in Game Theory}. Cambridge, MA, USA: MIT Press, 1994.

\bibitem{southey2005}
F. Southey, M. Bowling, B. Larson, C. Piccione, N. Burch, D. Billings, and C. Rayner, ``Bayes' bluff: Opponent modelling in poker,'' in \emph{Proc. 21st Conf. Uncertainty Artif. Intell. (UAI)}, 2005, pp. 550--558.

\bibitem{ganzfried2011}
S. Ganzfried and T. Sandholm, ``Game theory-based opponent modeling in large imperfect-information games,'' in \emph{Proc. 10th Int. Conf. Auton. Agents Multiagent Syst. (AAMAS)}, 2011, pp. 533--540.

\bibitem{johanson2007}
M. Johanson, M. Zinkevich, and M. Bowling, ``Computing robust counter-strategies,'' in \emph{Advances in Neural Information Processing Systems}, vol. 20, 2007, pp. 721--728.

\bibitem{ganzfried2015}
S. Ganzfried and T. Sandholm, ``Safe opponent exploitation,'' \emph{ACM Trans. Econ. Comput.}, vol. 3, no. 2, pp. 1--28, 2015.

\bibitem{silver2016}
D. Silver \emph{et al.}, ``Mastering the game of Go with deep neural networks and tree search,'' \emph{Nature}, vol. 529, no. 7587, pp. 484--489, 2016.

\bibitem{vinyals2019}
O. Vinyals \emph{et al.}, ``Grandmaster level in StarCraft II using multi-agent reinforcement learning,'' \emph{Nature}, vol. 575, no. 7782, pp. 350--354, 2019.

\bibitem{lanctot2017}
M. Lanctot, V. Zambaldi, A. Gruslys, A. Lazaridou, K. Tuyls, J. P\'erolat, D. Silver, and T. Graepel, ``A unified game-theoretic approach to multiagent reinforcement learning,'' in \emph{Advances in Neural Information Processing Systems}, vol. 30, 2017, pp. 4190--4203.

\bibitem{balduzzi2019}
D. Balduzzi, M. Garnelo, Y. Bachrach, W. M. Czarnecki, J. P\'erolat, M. Jaderberg, and T. Graepel, ``Open-ended learning in symmetric zero-sum games,'' in \emph{Proc. 36th Int. Conf. Mach. Learn. (ICML)}, 2019, pp. 434--443.

\bibitem{perezliebana2019}
D. Perez-Liebana, J. Liu, A. Khalifa, R. D. Gaina, J. Togelius, and S. M. Lucas, ``General video game AI: A multitrack framework for evaluating agents, games, and content generation algorithms,'' \emph{IEEE Trans. Games}, vol. 11, no. 3, pp. 195--214, 2019.

\bibitem{wellman2006}
M. P. Wellman, ``Methods for empirical game-theoretic analysis (extended abstract),'' in \emph{Proc. 21st Nat. Conf. Artif. Intell. (AAAI)}, 2006, pp. 1552--1556.

\bibitem{camerer2003}
C. F. Camerer, \emph{Behavioral Game Theory: Experiments in Strategic Interaction}. Princeton, NJ, USA: Princeton Univ. Press, 2003.

\bibitem{yannakakis2018}
G. N. Yannakakis and J. Togelius, \emph{Artificial Intelligence and Games}. Cham, Switzerland: Springer, 2018.

\bibitem{shannon1948}
C. E. Shannon, ``A mathematical theory of communication,'' \emph{Bell Syst. Tech. J.}, vol. 27, pp. 379--423, 623--656, 1948.

\bibitem{mookherjee1994}
D. Mookherjee and B. Sopher, ``Learning behavior in an experimental matching pennies game,'' \emph{Games Econ. Behav.}, vol. 7, no. 1, pp. 62--91, 1994.

\bibitem{zinkevich2007}
M. Zinkevich, M. Johanson, M. Bowling, and C. Piccione, ``Regret minimization in games with incomplete information,'' in \emph{Advances in Neural Information Processing Systems}, vol. 20, 2007, pp. 1729--1736.

\bibitem{bowling2015}
M. Bowling, N. Burch, M. Johanson, and O. Tammelin, ``Heads-up limit hold'em poker is solved,'' \emph{Science}, vol. 347, no. 6218, pp. 145--149, 2015.

\bibitem{brown2018}
N. Brown and T. Sandholm, ``Superhuman AI for heads-up no-limit poker: Libratus beats top professionals,'' \emph{Science}, vol. 359, no. 6374, pp. 418--424, 2018.

\bibitem{brown2019}
N. Brown and T. Sandholm, ``Superhuman AI for multiplayer poker,'' \emph{Science}, vol. 365, no. 6456, pp. 885--890, 2019.

\end{thebibliography}

\end{document}
